\section{Metodología de investigación aplicada}
El presente proyecto se desarrolló bajo un enfoque cualitativo, orientado a la comprensión e interpretación de conceptos relacionados con el patrón Modelo-Vista-Controlador (MVC), la arquitectura de Servicios Web RESTful y la comparación de frameworks backend como Spring Boot y Node.js. La implementación práctica se gestionó bajo la metodología ágil Scrum.

\subsection{Fase 1: Recolección y construcción del thesaurus}
Se utilizó Google Scholar, realizando consultas con palabras clave como: MVC, Servicios RESTful, Spring Boot, Node.js, Seguridad en aplicaciones web. El propósito fue reunir evidencia científica que respaldara el uso del patrón MVC, la arquitectura REST y la idoneidad de los frameworks backend seleccionados para la construcción de sistemas robustos y seguros.

\subsection{Fase 2: Análisis cualitativo de la evidencia (síntesis del thesaurus)}
El objetivo fue transformar la información documental en argumentos sólidos que justificaran la elección de MVC y el análisis comparativo entre Spring Boot y Node.js.
•	Se realizó una lectura crítica para identificar las ventajas del MVC en términos de separación de responsabilidades, mantenibilidad y testabilidad.
•	Se analizaron las características clave de Spring Boot y Node.js para determinar su aplicabilidad en diferentes contextos de proyecto.
•	Se consideraron las vulnerabilidades y riesgos comunes en aplicaciones web (como inyección SQL o XSS) para enfocar la implementación en la seguridad del backend.

\subsection{Fase 3: Diseño y desarrollo del sistema}
Transformar los criterios teóricos y metodológicos en un sistema funcional de servicios web, aplicando la arquitectura MVC y Servicios RESTful como guías de construcción y organización del proyecto.

\subsubsection{Procedimiento detallado}
\paragraph{1. Diseño de la Arquitectura (MVC con REST):}
o	Se definieron las capas principales: Modelo (lógica de negocio y persistencia de datos), y Controlador (gestión de las peticiones REST, mapeando verbos HTTP a operaciones del Modelo).
o	Esta separación permite mejorar la mantenibilidad y la escalabilidad del backend.

\paragraph{2. Selección y aplicación de Frameworks (Spring Boot como base):}
o	Se priorizó Spring Boot (Java) para la implementación base de los servicios RESTful, debido a su robustez y alto nivel de seguridad, ideal para un entorno formativo que simula exigencias empresariales.
o	Se utilizó el patrón Programación Orientada a Objetos (POO) en Java para estructurar el código bajo principios de herencia, encapsulamiento y polimorfismo, garantizando un desarrollo ordenado y reutilizable.

\paragraph{3. Metodología Ágil (Scrum):}
o	Se implementó la metodología Scrum, realizando reuniones clave (como la Daily Scrum) y trabajando por sprints definidos, lo que permitió ajustar prioridades y detectar errores de manera temprana.

\subsubsection{Herramientas utilizadas}
•	Backend: Java con Spring Boot (por su robustez y seguridad).
•	Base de datos: PostgreSQL (para gestionar la información de recursos).
•	Herramientas de Desarrollo: Visual Studio Code (VSC) (IDE ligero y flexible).
•	Gestión de Dependencias: Maven (para la gestión de librerías y ciclo de vida de proyectos Java).
•	Control de versiones: GitHub (para trabajo colaborativo y control de cambios).